% Generated draft: requires mathematical and source review.
\documentclass[11pt]{article}
\usepackage[margin=25mm]{geometry}
\usepackage{fontspec}
\setmainfont{Latin Modern Roman}
\usepackage{amsmath,amssymb,mathtools,unicode-math}
\setmathfont{Latin Modern Math}
\usepackage{enumitem,needspace,xurl,hyperref}
\setlength{\parindent}{0pt}
\setlength{\parskip}{5pt}
\setlength{\emergencystretch}{3em}
\newcommand{\N}{\mathbb{N}}
\newcommand{\Z}{\mathbb{Z}}
\newcommand{\Q}{\mathbb{Q}}
\newcommand{\R}{\mathbb{R}}
\newcommand{\C}{\mathbb{C}}
\newcommand{\F}{\mathbb{F}}
\newcommand{\A}{\mathbb{A}}
\newcommand{\PP}{\mathbb{P}}
\newcommand{\E}{\mathbb{E}}
\newcommand{\eps}{\varepsilon}
\newcommand{\dd}{\,\mathrm{d}}
\begin{document}
\section*{UnsolvedMath ID: 30002383}
\textbf{Source problem code:} catalogue entry
\subsection{Definitions and mathematical statement}
Let \(P\) be a finite set of rational primes satisfying
\[
\exists G\;\bigl(G\text{ is a finite nonabelian simple group and }\{p\text{ prime}:p\mid |G|\}\subseteq P\bigr).
\]
For each integer \(m\geq 1\), let \(S_m\) be the symmetric group on \(m\) letters and \(A_m\) its alternating subgroup. For a number field \(K/\Q\) of degree \(m\), write \(\widetilde K\) for its normal closure over \(\Q\). Interpret the source's phrase ``Galois group \(A_m\) or \(S_m\)'' as referring to \(\operatorname{Gal}(\widetilde K/\Q)\), acting on the \(m\) embeddings of \(K\) into \(\widetilde K\). Define
\[
F_P(m)=\#\left\{\begin{array}{l|l}
[K]_{\Q} & \begin{gathered}
[K:\Q]=m,\\
\operatorname{Gal}(\widetilde K/\Q)=A_m\text{ or }S_m
\text{ in its natural permutation action},\\
\text{every rational prime ramified in }K/\Q\text{ belongs to }P
\end{gathered}
\end{array}\right\},
\]
where \([K]_{\Q}\) denotes the \(\Q\)-isomorphism class of \(K\), and the permutation group is identified after labeling the embeddings. The ramification condition concerns finite primes; no condition at the real place is supplied.

\textbf{Unboundedness conjecture.} For every finite set \(P\) satisfying the displayed group-order condition, the sequence \((F_P(m))_{m\geq1}\) is unbounded. Explicitly,
\[
\forall B\in\Z_{\geq0}\;\exists m\in\Z_{\geq1}\qquad F_P(m)>B.
\]
The normal-closure interpretation and counting up to \(\Q\)-isomorphism are standard conventions made explicit here; the supplied record does not expressly state them. The source's report-year discrepancy and incomplete supporting literature attribution remain unresolved.
\subsection{Short English statement}
Fix any finite set \(P\) of primes containing every prime divisor of the order of at least one finite nonabelian simple group. The conjecture says that the number of degree-\(m\) number fields, counted up to \(\Q\)-isomorphism, whose normal closures have alternating or symmetric Galois group and which ramify only within \(P\), is unbounded as \(m\) varies. The supplied status is ``partially solved''; this draft does not independently verify that assessment.
\subsection{Sources}
\begin{itemize}
\item Local UnsolvedMath registry, \texttt{data/unsolvedmath/problems.json}: permanent ID \(30002383\), source problem code \texttt{catalogue entry}, title ``Unbounded Symmetric and Alternating Extensions with Restricted Ramification,'' in the Oberwolfach Reports Open Problems collection. This citation identifies the supplied registry record; no independent file or source inspection was performed.
\item Supplied source citation: \emph{Explicit Methods in Number Theory (2014)}, Oberwolfach Reports, pp.\ 2035--2082, DOI \texttt{10.4171/owr/2013/35}. Supplied source and literature link: \url{https://doi.org/10.4171/owr/2013/35}, labeled ``Oberwolfach Report 35/2013.'' The record gives proposed year 2013 and supplies no individual proposer.
\item Literature assessment dated exactly \texttt{2026-08-21}: supplied dated curation, not independently verified. The record labels the problem ``partially solved'' and states that Roberts's geometric fullness theorem and explicit Hurwitz-field examples support the conjecture, while the required arithmetic specialization statement proving unbounded \(F_P(m)\) remains open. No separate bibliographic references for those supporting claims are supplied. The background expressly requests independent verification of time-sensitive status and cited claims.
\end{itemize}
\end{document}
